\section{Cosmic Time and the World-Cycles}

Buddhist time has no edge. There is no first moment and no creator. Worlds form, last, are destroyed, and stay empty, then form again, without beginning and without end. This is the most distinctive feature of the Buddhist cosmos. It runs on causation and karma, not on a maker.

\subsection{Beginningless Time}

Samsara has \textbf{no first moment}. The technical term is \textbf{beginningless} (Sanskrit anadi, Tibetan thog ma med pa). The reasoning is causal. Every moment of mind arises from a previous moment of mind. Every effect has a cause. So trace the chain back and you never reach a starting point. A first moment would be a moment without a prior cause, which the system forbids.

This is why Buddhism posits \textbf{no creator}. The classical texts argue against a creator god (Ishvara) directly. A first cause is incoherent, since a cause must itself have a cause. A perfect, unchanging creator cannot be the source of a changing world, because acting is itself a change. Suffering and the moral order are already explained by karma, so no designer or judge is needed. The question of what came before the beginning is rejected as malformed, because there is no beginning.

\subsection{The Kalpa and Its Kinds}

Time is measured in \textbf{kalpas} (Tibetan bskal pa), vast eons. A kalpa is roughly the lifespan of a world-system, a stretch of time almost beyond imagining. The tradition counts several kinds.

\begin{longtable}{@{}
L{3.4cm}L{8.0cm}@{}}
\toprule
Kind & What it measures \\
\midrule
Intermediate (antarakalpa) & One rise-and-fall of human lifespan. The base unit. \\
Incalculable (asankhyeya) & An uncountable eon. A bodhisattva needs three of these to reach buddhahood. \\
Great (mahakalpa) & One full world-cycle. Equal to eighty intermediate kalpas. \\
\bottomrule
\end{longtable}

A common simile gives the scale. A mountain one yojana on each side is brushed once a century by a passing celestial's silk robe. The time to wear the mountain flat is less than one kalpa.

\subsection{The Four Phases of a World-Cycle}

A great kalpa runs through \textbf{four phases}, each lasting twenty intermediate kalpas, for eighty in all.

\begin{longtable}{@{}
rL{2.8cm}L{5.4cm}L{2.4cm}@{}}
\toprule
\# & Phase & Description & Length \\
\midrule
1 & Formation & The world re-forms from the wind-base up. Beings descend from the Abhasvara heaven. & 20 kalpas \\
2 & Duration & The world is stable and inhabited. Lifespan cycles happen here. & 20 kalpas \\
3 & Destruction & The world is destroyed by fire, water, or wind. Beings flee upward. & 20 kalpas \\
4 & Emptiness & Only empty space remains until the next formation. & 20 kalpas \\
\bottomrule
\end{longtable}

So the cycle is \textbf{formation, duration, destruction, emptiness}, then formation again, without end. We live now in the duration phase of the current great kalpa, on the declining arc, with human lifespan said to be falling from its last peak toward the next trough.

\subsection{The Three Destructions}

The world is destroyed at the end of each cycle, but not always the same way and not always to the same height. There are three kinds.

\begin{longtable}{@{}
L{2.4cm}L{6.6cm}L{2.2cm}@{}}
\toprule
Destruction & Spares everything from this level up & Frequency \\
\midrule
Fire & the Abhasvara heaven (2nd dhyana) and above & most common \\
Water & the Subhakrtsna heaven (3rd dhyana) and above & rarer \\
Wind & the Brhatphala heaven (4th dhyana) and above & rarest \\
\bottomrule
\end{longtable}

The destructions run on a schedule. Seven fires, then one water, repeated seven times, then one wind. One full super-cycle is forty-nine fires, seven waters, and one wind, sixty-four destructions in all. The fourth-dhyana plane and the Pure Abodes above it survive all three, but they too end when their beings' karma runs out.

\subsection{The Fortunate Aeon}

The present great kalpa is a \textbf{fortunate aeon} (\textbf{Bhadrakalpa}). It earns the name because \textbf{one thousand Buddhas} are prophesied to appear within it. Shakyamuni, the historical Buddha, was the fourth. \textbf{Maitreya} will be the fifth, descending from the Tusita heaven when the time is right. The vast majority of the thousand are still to come.

\subsection{The Kalachakra Mapping}

The \textbf{Wheel of Time} (\textbf{Kalachakra}, Tibetan dus kyi 'khor lo) is a major tantric system that maps cosmic time onto the body. It works across three levels. The \textbf{outer Kalachakra} is the cosmos and its cycles, the motions of planets and stars. The \textbf{inner Kalachakra} is the body, its channels and winds, and the cycles of the breath. The \textbf{other Kalachakra} is the meditation practice that aligns the two. The system ties the great cycles of cosmic time to the cycles of breath and life, so that the turning of the universe and the turning of the breath are read as one rhythm. Like the rest of the Buddhist view, it presents time as cyclic and beginningless.
